\Askhsh{21121_SOLUTION}{1}
\begin{ylh}{1}
    \item [Ύλη από εικόνα]
\end{ylh}
\begin{alist}
\item Επειδή $ΒΓ = 13$ και $ΑΔ = 6$, επομένως το εμβαδόν του τριγώνου $ΑΒΓ$ είναι
\[
(ΑΒΓ) = \frac{1}{2} \cdot ΒΓ \cdot ΑΔ \text{ ή } (ΑΒΓ) = \frac{1}{2} \cdot 13 \cdot 6 \text{ ή } (ΑΒΓ) = 39.
\]

\item
\begin{rlist}
\item Ο κυκλικός τομέας $Α\overset{\frown}{ΕΔΖ}$ είναι γωνίας $\mu = \hat{Α} = 90^\circ$ και ακτίνας $R = ΑΔ = 6$, οπότε το εμβαδόν του είναι
\[
(Α\overset{\frown}{ΕΔΖ}) = \frac{\pi R^2 \mu}{360^\circ} = \frac{\pi \cdot 6^2 \cdot 90^\circ}{360^\circ} = 9\pi.
\]
\item Το εμβαδόν $E$ του χωρίου που είναι εσωτερικά του τριγώνου $ΑΒΓ$ και εξωτερικά του κύκλου είναι ίσο με τη διαφορά των εμβαδών του τριγώνου $ΑΒΓ$ και του κυκλικού τομέα $Α\overset{\frown}{ΕΔΖ}$. Επομένως είναι
\[
E = (ΑΒΓ) - (Α\overset{\frown}{ΕΔΖ}) = 39 - 9\pi.
\]
\end{rlist}
\end{alist}